\subsection{Soldering Iron}
\markright{Soldering Iron}

\subsubsection*{Definition}
A soldering iron is a hand-held tool with an electrically heated metal tip (the
\textit{bit}) used to melt solder and join component leads to printed circuit board
pads or wire terminals \cite{wiki:solderiron}. Solder is a low-melting-point metal
alloy that, when melted and solidified, forms a mechanically strong and electrically
conductive joint. Traditional tin--lead solder (60Sn/40Pb) melts at approximately
183\,$^\circ$C; modern lead-free alloys (e.g., Sn96.5/Ag3.0/Cu0.5, ``SAC305'') melt at
217--227\,$^\circ$C \cite{wiki:solderiron}.

\labimage{lab-01/assets/soldering_iron.jpeg}%
  {A general-purpose soldering iron ($\approx$30--40\,W) of the type used
   in this laboratory; the replaceable tip (bit) is heated resistively to the
   set temperature.}{solderiron_lab}

\subsubsection*{Specifications}
The laboratory soldering iron is rated approximately 30--40\,W, operating at mains
voltage (220\,V AC). Tip temperature for general electronics work with lead-free solder
is set to 315--370\,$^\circ$C. Temperature-controlled soldering stations maintain the
set temperature to within a few degrees via a closed-loop thermocouple sensor, reducing
the risk of thermal damage to heat-sensitive components \cite{wiki:solderiron}.

\subsubsection*{Purpose of Use}
The soldering iron is used to create permanent solder joints on veroboard and PCBs,
to tin (pre-coat) wire ends and component leads, and to desolder components for rework.
A good solder joint is smooth, shiny, and conical; a defective ``cold joint'' is dull
and grainy, with unreliable electrical conductivity.

\subsubsection*{Applications}
Soldering irons are used in PCB assembly and repair, wire-to-terminal connections,
cable harness manufacturing, educational laboratory circuit assembly, and field-service
rework of electronic equipment \cite{wiki:solderiron}.
